\documentclass[11pt,a4paper]{article}
\usepackage[utf8]{inputenc}
\usepackage{amsmath}
\usepackage{graphicx}
\usepackage{hyperref}
\usepackage{geometry}
\geometry{margin=1in}

\title{\textbf{Reinforcement-Learned Adaptive Prosody for Text-to-Speech Systems}}
\author{Arush Singhal}
\date{March 2026}

\begin{document}

\maketitle

\begin{abstract}
Current Text-to-Speech (TTS) models successfully generate intelligible speech but struggle with diverse, expressive prosody when trained purely via reconstruction loss. We present a simulated lightweight framework demonstrating how Reinforcement Learning (RL) can explicitly optimize pitch, duration, energy, and pausing. Through a custom continuous-action environment and a composite proxy reward, we reveal critical tradeoffs between naturalness and intelligibility.
\end{abstract}

\section{Introduction}
Prosody encompasses the rhythm, stress, and intonation of speech. It is highly subjective, creating a one-to-many mapping problem in standard TTS training frameworks. In this note, we propose an RL perspective, allowing TTS models to explore the prosodic landscape rather than converging on overly smoothed regression averages.

\section{Related Motivation / Background}
Traditional systems employ FastSpeech or Tacotron architectures using Mean Squared Error (MSE) to predict pitch and duration. However, MSE heavily penalizes deviation from the single ground-truth recording, ignoring that alternative prosodic styles are equally valid. 

\section{Method}
We frame prosody selection as a continuous control problem. Let states $s_t$ represent text features and actions $a_t \in \mathbb{R}^4$ represent changes in pitch, duration constraint, energy, and pause delay. The policy $\pi_\theta(a|s)$ is optimized against a composite reward:
\begin{equation}
R = \lambda_1 R_{nat} + \lambda_2 R_{int} + \lambda_3 R_{ctx} + \lambda_4 R_{mono} - \lambda_5 R_{lat}
\end{equation}

\section{Experimental Setup}
We employ a sequence-based Hill Climbing optimization strategy (acting as a zero-order policy gradient stand-in) across four distinct scenarios:
1. Reward ablation on final policy parameters.
2. Sweep across $\lambda_{nat}$ vs $\lambda_{int}$ to find the Pareto frontier.
3. Testing contextual invariance on varied utterances.
4. Latency scaling.

\section{Results}
Results generated in \texttt{figures/} highlight quantitative differences. The agent successfully learned to insert appropriate pauses before commas and raise pitch at the end of questions solely through reward direction. 

\section{Discussion}
The key insight is the inherent tension between intelligibility (which prefers strict timings) and naturalness (which relies on variance). By deploying an RL methodology, TTS designers have direct levers to tune these traits for distinct applications (e.g. fast navigation voices vs. expressive audiobooks).

\section{Limitations}
This project is a rapid prototype simulation. It employs heuristic functions as surrogate human evaluators and assumes independent manipulatability of acoustic features without mapping artifacts.

\section{Conclusion}
Treating prosody as a control problem yields immense flexibility. RL provides a principled mechanism for addressing the subjective nature of human speech generation, escaping the "average-case" trap of supervised learning.

\end{document}
